% 图契约：电感电流增加时感生电场、磁场与向内能流的方向；几何、方向和公式与正文同一约定。
\documentclass[tikz,border=5pt]{standalone}
\usepackage{fontspec,xeCJK,amsmath,amssymb}
\IfFontExistsTF{Times New Roman}{\setmainfont{Times New Roman}}{\setmainfont{TeX Gyre Termes}}
\IfFontExistsTF{Heiti SC}{\setCJKmainfont[AutoFakeBold=1.5]{Heiti SC}}{\setCJKmainfont[AutoFakeBold=1.5]{FandolHei-Regular}}
\usetikzlibrary{arrows.meta,calc,patterns}
\definecolor{ink}{HTML}{243745}
\definecolor{blue}{HTML}{275DAB}
\definecolor{green}{HTML}{177C63}
\definecolor{red}{HTML}{BD4935}
\definecolor{orange}{HTML}{AA6C16}
\definecolor{purple}{HTML}{715393}
\begin{document}
\begin{tikzpicture}[x=1cm,y=1cm,font=\fontsize{9.5}{12}\selectfont,
v/.style={-{Latex[length=2mm,width=1.3mm]},line width=.85pt},
axis/.style={-{Latex[length=1.7mm,width=1.1mm]},ink!65,line width=.55pt},
aux/.style={ink!45,dashed,line width=.55pt},
lab/.style={fill=white,inner sep=1.8pt},
card/.style={draw=ink!25,fill=blue!3,rounded corners=2pt,text width=6.1cm,minimum height=.85cm,inner sep=6pt,align=center}]
\path[use as bounding box](-8,-3.6)rectangle(8,3.65);
\node[font=\bfseries\fontsize{10.5}{13}\selectfont,text width=15.1cm,align=center]at(0.00000,3.22000){M7-H07　电感电流增加时感生电场、磁场与向内能流的方向};
\begin{scope}[shift={(-4,0)}]
\draw[blue,fill=blue!5](0,0)circle(1.6);
\draw[ink!60,dashed](0,0)circle(2.25);
\draw[ink!50,dashed](0,0)circle(0.9);
\node[]at(-1.40000,1.30000){$R$};
\node[]at(1.80000,-1.80000){$r_>$};
\node[lab]at(0.65000,0.95000){$r_<$};
\draw[blue](0,0)circle(0.13);
\fill[blue](0,0)circle(.035);
\draw[blue](0.5,0.35)circle(0.13);
\fill[blue](0.5,0.35)circle(.035);
\draw[blue](-0.5,0.35)circle(0.13);
\fill[blue](-0.5,0.35)circle(.035);
\draw[blue](0.5,-0.35)circle(0.13);
\fill[blue](0.5,-0.35)circle(.035);
\draw[blue](-0.5,-0.35)circle(0.13);
\fill[blue](-0.5,-0.35)circle(.035);
\draw[v,red](.9,0)arc(0:-230:.9);
\draw[v,red](2.25,0)arc(0:-230:2.25);
\node[]at(0.00000,2.65000){$\dot B_z>0$};
\node[text=red]at(0.00000,-2.65000){$\mathbf E_{\rm ind}$ 顺时针；$\mathbf S$ 向内};
\draw[v,green](2.70000,0.00000)--(1.83600,0.00000);
\draw[v,green](-2.70000,0.00000)--(-1.83600,0.00000);
\end{scope}
\node[card]at(4.00000,1.25000){$\mathbf E_{\rm ind}\cdot\mathbf J<0$};
\node[card]at(4.00000,-1.05000){$P_{\rm in}=\frac d{dt}(\frac12LI^2)$};
\end{tikzpicture}\end{document}
